\documentclass[11pt,a4paper]{article}

\usepackage[utf8]{inputenc}
\usepackage[T1]{fontenc}
\usepackage[margin=2.5cm]{geometry}
\usepackage{hyperref}
\usepackage{xcolor}
\usepackage{listings}
\usepackage{enumitem}
\usepackage{booktabs}
\usepackage{longtable}
\usepackage{amsmath,amssymb}
\usepackage{fancyhdr}
\usepackage{titlesec}

\definecolor{isls-blue}{HTML}{1a5276}
\definecolor{isls-orange}{HTML}{e67e22}
\definecolor{isls-green}{HTML}{27ae60}
\definecolor{isls-red}{HTML}{c0392b}
\definecolor{isls-gray}{HTML}{7f8c8d}
\definecolor{code-bg}{HTML}{f7f9fb}

\lstdefinestyle{rust}{
  language=C,
  basicstyle=\ttfamily\small,
  keywordstyle=\color{isls-blue}\bfseries,
  commentstyle=\color{isls-gray}\itshape,
  stringstyle=\color{isls-green},
  backgroundcolor=\color{code-bg},
  frame=single,
  rulecolor=\color{isls-gray!30},
  breaklines=true,
  showstringspaces=false,
  tabsize=4,
  morekeywords={fn,let,mut,pub,struct,impl,enum,use,mod,self,Self,
    crate,super,where,trait,for,in,if,else,match,return,async,await,
    Ok,Err,Some,None,Vec,String,HashMap,HashSet,PathBuf,Result,Option},
}
\lstdefinestyle{bash}{
  language=bash,
  basicstyle=\ttfamily\small,
  backgroundcolor=\color{code-bg},
  frame=single,
  rulecolor=\color{isls-gray!30},
  breaklines=true,
  showstringspaces=false,
}
\lstdefinestyle{json}{
  basicstyle=\ttfamily\small,
  backgroundcolor=\color{code-bg},
  frame=single,
  rulecolor=\color{isls-gray!30},
  breaklines=true,
  showstringspaces=false,
}

\pagestyle{fancy}
\fancyhf{}
\fancyhead[L]{\small\textcolor{isls-gray}{ISLS D7 Spec --- Generationsspirale}}
\fancyhead[R]{\small\textcolor{isls-gray}{v1.0.0 --- \today}}
\fancyfoot[C]{\thepage}

\titleformat{\section}
  {\Large\bfseries\color{isls-blue}}{\thesection}{1em}{}
\titleformat{\subsection}
  {\large\bfseries\color{isls-blue!80}}{\thesubsection}{1em}{}
\titleformat{\subsubsection}
  {\normalsize\bfseries\color{isls-blue!60}}{\thesubsubsection}{1em}{}

\hypersetup{
  colorlinks=true,
  linkcolor=isls-blue,
  urlcolor=isls-orange,
  citecolor=isls-green,
}

\begin{document}

\begin{titlepage}
\centering
\vspace*{2.5cm}

{\Huge\bfseries\color{isls-blue} ISLS D7 Specification}\\[0.5cm]
{\LARGE\color{isls-orange} Generationsspirale: Das Cockpit}\\[0.3cm]
{\large\color{isls-gray} Architekt-Modus + Drucker-Modus + Messbare Verbesserung}\\[2cm]

{\large
\begin{tabular}{rl}
\textbf{System:}    & ISLS --- Intelligent Semantic Ledger Substrate \\
\textbf{Version:}   & v4.1 (D7) \\
\textbf{Author:}    & Sebastian Klemm \\
\textbf{Date:}      & \today \\
\textbf{Status:}    & Specification \\
\textbf{Depends:}   & D6 (complete), isls-gateway, isls-chat \\
\textbf{Branch:}    & \texttt{claude/implement-isls-d7-spec} \\
\end{tabular}
}

\vfill

{\small\color{isls-gray}
Mithras Substructure University\\
Repository: \texttt{https://github.com/LashSesh/graphity}\\
License: MIT
}
\end{titlepage}

\tableofcontents
\newpage


% ══════════════════════════════════════════════════════════════════════════════
\section{Executive Summary}
% ══════════════════════════════════════════════════════════════════════════════

D1--D6 built the engine. D7 builds the cockpit.

Until now, ISLS is operated through the CLI: \texttt{isls forge-chat 
-m "..."} --- single-shot, one sentence, one app. This is powerful 
for demos and testing, but inadequate for real work. Real projects 
need iteration: ``Add a rating system.'' ``Make the booking entity 
track cancellations.'' ``Actually, remove the frontend --- this should 
be a headless API.''

D7 transforms the existing \texttt{isls-gateway} into a two-mode 
cockpit:

\begin{description}[nosep]
\item[Architekt-Modus] Multi-turn conversation. You describe, the 
  system suggests entities and norms, you refine. A readiness 
  indicator shows what is defined and what is missing. The 
  conversation accumulates into an \texttt{AppSpec} incrementally.
\item[Drucker-Modus] One button: ``Forge.'' The accumulated 
  AppSpec flows through the HDAG pipeline (S0--S7). Real-time 
  progress via WebSocket. Result page with file browser.
\end{description}

The \emph{Generationsspirale} is the measurement: apps generated 
through the Cockpit (multi-turn, norm-enriched) vs. apps generated 
through the CLI (single-shot). D7 logs metrics for both paths so 
the comparison emerges naturally.

\subsection{What Exists Already}

The \texttt{isls-gateway} crate contains substantial infrastructure:

\begin{itemize}[nosep]
\item Axum HTTP server with 30+ routes.
\item Embedded Studio HTML (single-page app, dark theme, 9 views).
\item \texttt{POST /api/chat}: norm-aware intent extraction, plan 
      composition, pending plan storage.
\item \texttt{POST /api/projects/forge}: takes a plan ID, runs the 
      full forge pipeline, returns results.
\item WebSocket event hub (\texttt{/events}) with subscribe/publish.
\item In-memory project store, norm registry integration.
\item \texttt{POST /agent/chat}: background job runner for code 
      modifications.
\end{itemize}

D7 does not rewrite this. D7 extends it with: (1) multi-turn 
conversation state, (2) readiness check, (3) forge progress 
streaming, (4) Studio HTML upgrade.

\subsection{What D7 Achieves}

\begin{enumerate}[nosep]
\item \textbf{Conversation Sessions:} Multi-turn chat that 
      accumulates entities, relationships, and constraints.
\item \textbf{Readiness Check:} API that evaluates the current 
      session state and returns a completeness assessment.
\item \textbf{Live Forge:} Forge execution with WebSocket progress 
      events (node-by-node, compile status, final result).
\item \textbf{Studio Upgrade:} Chat view, readiness indicator, 
      forge button, progress view, result file browser.
\item \textbf{Generation Metrics:} Each generation logs token 
      count, compile success, entity count, duration --- enabling 
      CLI vs. Cockpit comparison.
\end{enumerate}

\subsection{What D7 is NOT}

\begin{itemize}[nosep]
\item \textbf{Not a rewrite of isls-gateway.} Existing routes 
      and handlers remain. New endpoints are added alongside.
\item \textbf{Not a full IDE.} No code editor, no file tree 
      editing, no git integration. The Studio shows results, 
      it doesn't edit them.
\item \textbf{Not multi-user.} Single operator. No accounts 
      beyond what the generated app provides.
\item \textbf{Not persistent storage.} Sessions and projects 
      are in-memory for now. Persistence is D8.
\end{itemize}


% ══════════════════════════════════════════════════════════════════════════════
\section{The Two Modes}
% ══════════════════════════════════════════════════════════════════════════════

\subsection{Architekt-Modus: Conversation $\to$ AppSpec}

The architect conversation is a sequence of messages that 
incrementally builds an \texttt{AppSpec}:

\begin{lstlisting}[style=bash]
USER: "I need a system to manage my crypto trading journal.
       Trades, portfolios, and performance metrics."

ISLS: Found 3 entities: Trade, Portfolio, PerformanceMetric.
      Activated norms: CRUD-Entity, JWT-Auth, Dashboard-KPI.
      Infrastructure: Web-Server + Database + Auth + Frontend.
      
      Missing: Trade has no fields defined yet. What should
      a Trade look like?

USER: "Each trade has: pair (like BTC/USDT), side (buy/sell), 
       entry price, exit price, quantity, timestamp, notes."

ISLS: Updated Trade with 7 fields. Added FK: Trade -> Portfolio.
      Readiness: 72% (PerformanceMetric has no fields yet).

USER: "Performance metrics are computed — average win rate, 
       total PnL, Sharpe ratio. Actually, make this a 
       headless API, no frontend."

ISLS: Updated PerformanceMetric (3 computed fields).
      Removed Frontend infrastructure.
      Readiness: 95% — ready to forge.
      
      [FORGE] button active.
\end{lstlisting}

Each user message is sent via \texttt{POST /api/session/\{id\}/message}.
The backend:
\begin{enumerate}[nosep]
\item Sends the conversation history + current AppSpec to the LLM.
\item LLM returns a structured JSON with entity updates, new 
      entities, removed entities, infrastructure changes.
\item Backend merges the updates into the session's AppSpec.
\item Returns the updated state + readiness assessment.
\end{enumerate}

\subsection{Drucker-Modus: AppSpec $\to$ Compiled Application}

When the readiness check passes (or the user forces it), the 
``Forge'' button triggers:

\begin{enumerate}[nosep]
\item Convert session AppSpec to TOML (using existing 
      \texttt{json\_to\_toml()}).
\item Write TOML to temp file.
\item Call \texttt{forge-v2} pipeline with the TOML.
\item Stream progress events via WebSocket:
      \begin{itemize}[nosep]
      \item \texttt{forge:node} --- for each HDAG node processed.
      \item \texttt{forge:compile} --- cargo check started/passed/failed.
      \item \texttt{forge:coagula} --- S6 fix cycle (if triggered).
      \item \texttt{forge:complete} --- final result with stats.
      \end{itemize}
\item Store result in project store.
\item Display result page with file browser.
\end{enumerate}


% ══════════════════════════════════════════════════════════════════════════════
\section{Work Packages}
% ══════════════════════════════════════════════════════════════════════════════

% ──────────────────────────────────────────────────────────────────────────────
\subsection{W1: Conversation Sessions}
\label{sec:w1}
% ──────────────────────────────────────────────────────────────────────────────

\subsubsection{Goal}

Add session-based multi-turn conversation state to the gateway.

\subsubsection{Data Structures}

\begin{lstlisting}[style=rust]
/// A conversation session accumulating toward an AppSpec.
pub struct ArchitectSession {
    pub id: String,
    pub created_at: String,
    /// Full conversation history.
    pub messages: Vec<SessionMessage>,
    /// The evolving AppSpec (starts empty, grows per turn).
    pub spec: AppSpec,
    /// Activated infrastructure norms.
    pub infra_norms: Vec<String>,
    /// Blueprint derived from infra norms.
    pub blueprint: InfraBlueprint,
    /// API key for LLM calls (provided at session creation).
    pub api_key: Option<String>,
    /// LLM model name.
    pub model: String,
    /// Generation result (after forge).
    pub forge_result: Option<ForgeResult>,
}

pub struct SessionMessage {
    pub role: String,    // "user" or "assistant"
    pub content: String,
    pub timestamp: String,
    /// Structured updates extracted from this message.
    pub updates: Option<SpecUpdate>,
}

pub struct SpecUpdate {
    pub entities_added: Vec<String>,
    pub entities_modified: Vec<String>,
    pub entities_removed: Vec<String>,
    pub infra_changes: Vec<String>,
}
\end{lstlisting}

\subsubsection{New Endpoints}

\begin{longtable}{llp{6cm}}
\toprule
\textbf{Method} & \textbf{Path} & \textbf{Purpose} \\
\midrule
POST & \texttt{/api/session} 
  & Create new session. Body: \texttt{\{api\_key, model?\}}. 
    Returns session ID. \\
GET & \texttt{/api/session/\{id\}} 
  & Get session state: messages, current AppSpec, readiness. \\
POST & \texttt{/api/session/\{id\}/message} 
  & Send user message. Triggers LLM call. Returns assistant 
    response + updated AppSpec + readiness. \\
DELETE & \texttt{/api/session/\{id\}} 
  & Delete session. \\
GET & \texttt{/api/sessions} 
  & List all active sessions. \\
\bottomrule
\end{longtable}

\subsubsection{LLM Conversation Protocol}

When \texttt{/api/session/\{id\}/message} is called:

\begin{enumerate}[nosep]
\item Build LLM prompt:
      \begin{itemize}[nosep]
      \item System message: ``You are an architecture assistant for 
            ISLS. Extract entities, fields, relationships, and 
            infrastructure requirements from the conversation. 
            Return a JSON object with updates.''
      \item Conversation history (all previous messages).
      \item Current AppSpec as context.
      \item User's new message.
      \end{itemize}
\item Call Oracle (OpenAI API) with the prompt.
\item Parse response: the LLM returns both a human-readable 
      message AND a structured JSON block (fenced with 
      \texttt{```json ... ```}).
\item Apply the structured updates to the session's AppSpec.
\item Recompute readiness.
\item Return the assistant message + updated state.
\end{enumerate}

The LLM prompt instructs the model to return a specific JSON schema:

\begin{lstlisting}[style=json]
{
  "message": "Human-readable response to the user.",
  "updates": {
    "entities_add": [
      { "name": "Trade", "fields": [
        { "name": "pair", "field_type": "String" },
        { "name": "entry_price", "field_type": "f64" }
      ], "foreign_keys": [] }
    ],
    "entities_modify": [
      { "name": "Portfolio", "add_fields": [
        { "name": "total_value", "field_type": "f64" }
      ], "remove_fields": [] }
    ],
    "entities_remove": [],
    "infra": {
      "remove_frontend": true,
      "add_cli": false
    }
  }
}
\end{lstlisting}

\subsubsection{Fallback: No LLM Available}

If no API key is provided or the call fails, the session still 
works in manual mode: the user can directly add/modify entities 
via \texttt{POST /api/session/\{id\}/entity}. The architect 
conversation degrades gracefully to a structured form.

\subsubsection{New Files}

\begin{longtable}{lp{10cm}}
\toprule
\textbf{File} & \textbf{Purpose} \\
\midrule
\texttt{isls-gateway/src/session.rs}
  & \texttt{ArchitectSession}, \texttt{SessionMessage}, 
    \texttt{SpecUpdate}, session management logic. \\
\texttt{isls-gateway/src/architect.rs}
  & LLM conversation protocol: prompt building, response 
    parsing, AppSpec merging. \\
\bottomrule
\end{longtable}

\subsubsection{Modified Files}

\begin{longtable}{lp{10cm}}
\toprule
\textbf{File} & \textbf{Change} \\
\midrule
\texttt{isls-gateway/src/lib.rs}
  & Add session routes, session store to AppState, 
    \texttt{pub mod session; pub mod architect;} \\
\bottomrule
\end{longtable}

\subsubsection{W1 Verification}

\begin{enumerate}[nosep]
\item Create session: \texttt{POST /api/session} $\to$ 201 with 
      session ID.
\item Send message: \texttt{POST /api/session/\{id\}/message} 
      with ``Build me a pet shop app'' $\to$ response with 
      entities and readiness.
\item Send follow-up: ``Add a vaccination record for each pet'' 
      $\to$ entity modified, readiness updated.
\item Get state: \texttt{GET /api/session/\{id\}} shows full 
      conversation + accumulated AppSpec.
\end{enumerate}


% ──────────────────────────────────────────────────────────────────────────────
\subsection{W2: Readiness Check}
\label{sec:w2}
% ──────────────────────────────────────────────────────────────────────────────

\subsubsection{Goal}

Evaluate the current session state and report what is complete, 
what is missing, and whether the AppSpec is ready for forge.

\subsubsection{Readiness Criteria}

\begin{longtable}{lp{8cm}l}
\toprule
\textbf{Criterion} & \textbf{Description} & \textbf{Weight} \\
\midrule
\texttt{has\_entities}
  & At least 1 non-User entity exists.
  & Required \\
\texttt{entities\_have\_fields}
  & Every entity has $\geq$1 field.
  & Required \\
\texttt{has\_app\_name}
  & AppSpec has a non-empty app\_name.
  & Required \\
\texttt{fk\_targets\_valid}
  & All FK targets reference existing entities.
  & Required \\
\texttt{has\_description}
  & AppSpec has a non-empty description.
  & +10\% \\
\texttt{entities\_have\_types}
  & All fields have valid types (String, i32, i64, f64, bool).
  & +15\% \\
\texttt{has\_relationships}
  & At least 1 FK exists.
  & +10\% \\
\texttt{multi\_entity}
  & $\geq$3 entities (richer domain).
  & +15\% \\
\bottomrule
\end{longtable}

``Required'' criteria must all pass for the forge button to activate.
Percentage criteria contribute to the readiness score (0--100\%).

\subsubsection{API}

\begin{lstlisting}[style=json]
GET /api/session/{id}/readiness

Response:
{
  "ready": true,
  "score": 85,
  "criteria": {
    "has_entities": { "met": true, "detail": "3 entities" },
    "entities_have_fields": { "met": true },
    "has_app_name": { "met": true, "detail": "crypto-journal" },
    "fk_targets_valid": { "met": true },
    "has_description": { "met": true, "score": 10 },
    "entities_have_types": { "met": false, 
      "detail": "PerformanceMetric.sharpe_ratio has no type",
      "score": 0 },
    "has_relationships": { "met": true, "score": 10 },
    "multi_entity": { "met": true, "score": 15 }
  },
  "suggestions": [
    "Add type to PerformanceMetric.sharpe_ratio (String/i32/f64/bool)"
  ]
}
\end{lstlisting}

The readiness check is also computed automatically after each 
\texttt{/api/session/\{id\}/message} response and included in the 
response body.

\subsubsection{W2 Verification}

\begin{enumerate}[nosep]
\item Empty session: readiness = 0, \texttt{ready = false}.
\item Session with 1 entity + fields: readiness $>$0, 
      \texttt{ready = true} (required criteria met).
\item Session with incomplete fields: \texttt{suggestions} 
      lists the missing items.
\end{enumerate}


% ──────────────────────────────────────────────────────────────────────────────
\subsection{W3: Live Forge Integration}
\label{sec:w3}
% ──────────────────────────────────────────────────────────────────────────────

\subsubsection{Goal}

Connect the session's accumulated AppSpec to the forge pipeline 
with real-time WebSocket progress streaming.

\subsubsection{Endpoint}

\begin{lstlisting}[style=bash]
POST /api/session/{id}/forge
Body: { "output_dir": "./output/my-project" }  (optional)
\end{lstlisting}

\subsubsection{Pipeline}

\begin{enumerate}[nosep]
\item Validate readiness (required criteria must pass).
\item Convert AppSpec to TOML string via \texttt{json\_to\_toml()}.
\item Write TOML to temporary file.
\item Derive InfraBlueprint from session's infra norms.
\item Build ForgePlan (spec + norm\_ids + blueprint).
\item Spawn background task running \texttt{StagedClosure::execute()}.
\item During execution, publish WebSocket events:
\end{enumerate}

\subsubsection{WebSocket Event Types}

\begin{lstlisting}[style=json]
{ "type": "forge:start", "session_id": "...", 
  "total_nodes": 50 }

{ "type": "forge:node", "index": 21, "total": 50,
  "path": "backend/src/database/norm_queries.rs",
  "method": "llm", "tokens": 2340 }

{ "type": "forge:compile", "status": "running" }
{ "type": "forge:compile", "status": "passed" }
{ "type": "forge:compile", "status": "failed", 
  "errors": 12 }

{ "type": "forge:coagula", "cycle": 1, "max": 3 }

{ "type": "forge:complete", "success": true,
  "files": 50, "loc": 2535, "tokens": 17263,
  "duration_secs": 232.9, "output_dir": "..." }
\end{lstlisting}

\subsubsection{Implementation}

The forge runs in a \texttt{tokio::spawn} background task.
The challenge: \texttt{StagedClosure} is synchronous 
(blocking \texttt{cargo check}). Solution: wrap the forge 
execution in \texttt{tokio::task::spawn\_blocking()} and use 
a channel (\texttt{tokio::sync::mpsc}) to send progress events 
back to the async context, which publishes them to the WebSocket 
hub.

The \texttt{StagedClosure} needs minor modification: a callback 
hook after each node is processed (currently it only prints to 
stderr). Add an optional \texttt{progress\_callback: Option<Box<dyn 
Fn(ForgeProgress) + Send>>} to \texttt{StagedClosure} that fires 
after each node.

\subsubsection{Modified Files}

\begin{longtable}{lp{10cm}}
\toprule
\textbf{File} & \textbf{Change} \\
\midrule
\texttt{isls-gateway/src/lib.rs}
  & Add \texttt{POST /api/session/\{id\}/forge} route. \\
\texttt{isls-gateway/src/session.rs}
  & Add \texttt{forge\_session()} that converts session state 
    to ForgePlan and runs pipeline. \\
\texttt{isls-forge-llm/src/staged\_closure.rs}
  & Add optional \texttt{progress\_callback} field. Fire 
    callback after each node in S4 Solve. \\
\bottomrule
\end{longtable}

\subsubsection{W3 Verification}

\begin{enumerate}[nosep]
\item Create session, send messages until ready.
\item \texttt{POST /api/session/\{id\}/forge} --- starts 
      background forge.
\item WebSocket client receives \texttt{forge:start}, 
      \texttt{forge:node} (50x), \texttt{forge:compile}, 
      \texttt{forge:complete} events.
\item Session's \texttt{forge\_result} is populated.
\item Output directory contains compiled application.
\end{enumerate}


% ──────────────────────────────────────────────────────────────────────────────
\subsection{W4: Generation Metrics}
\label{sec:w4}
% ──────────────────────────────────────────────────────────────────────────────

\subsubsection{Goal}

Log metrics for every generation (CLI and Cockpit) to enable 
the Gen~2 vs.\ Gen~1 comparison.

\subsubsection{Metrics Structure}

\begin{lstlisting}[style=rust]
pub struct GenerationMetrics {
    pub id: String,
    pub timestamp: String,
    pub source: GenerationSource, // Cli | Cockpit
    pub description: String,
    pub entity_count: usize,
    pub file_count: usize,
    pub structural_files: usize,
    pub llm_files: usize,
    pub total_tokens: u64,
    pub compile_success: bool,
    pub coagula_cycles: u32,
    pub duration_secs: f64,
    pub norms_activated: Vec<String>,
    pub conversation_turns: usize, // 1 for CLI, N for Cockpit
}

pub enum GenerationSource {
    Cli,      // forge-chat, forge-v2
    Cockpit,  // /api/session/{id}/forge
}
\end{lstlisting}

\subsubsection{Storage}

Metrics are appended to \texttt{\textasciitilde/.isls/metrics.jsonl} 
(one JSON object per line, append-only). Both CLI 
(\texttt{staged\_closure.rs}) and gateway (\texttt{session.rs}) 
write to the same file.

\subsubsection{CLI Inspection}

\begin{lstlisting}[style=bash]
isls metrics                    # Summary
isls metrics --compare          # CLI vs Cockpit comparison
isls metrics --last 10          # Last 10 generations
\end{lstlisting}

\subsubsection{The Spiral Measurement}

After enough generations through both paths, \texttt{isls metrics 
--compare} shows:

\begin{lstlisting}[style=bash]
ISLS Generation Metrics — CLI vs Cockpit
---
                    CLI (Gen 1)    Cockpit (Gen 2)
Generations:        12             5
Avg entities:       5.3            7.8
Avg tokens:         12,400         15,100
Compile success:    83%            100%
Avg coagula:        0.42           0.00
Avg duration:       180s           240s
Conversation turns: 1.0            4.2
\end{lstlisting}

If Cockpit apps compile more reliably and have richer entity 
models (due to multi-turn refinement), Gen~2 $>$ Gen~1 is proven.

\subsubsection{New Files}

\begin{longtable}{lp{10cm}}
\toprule
\textbf{File} & \textbf{Purpose} \\
\midrule
\texttt{isls-forge-llm/src/metrics.rs}
  & \texttt{GenerationMetrics} struct, \texttt{append\_metrics()},
    \texttt{load\_metrics()}, comparison logic. \\
\texttt{isls-cli/src/cmd\_metrics.rs}
  & \texttt{isls metrics} subcommand handlers. \\
\bottomrule
\end{longtable}

\subsubsection{W4 Verification}

\begin{enumerate}[nosep]
\item Run \texttt{isls forge-chat} --- verify 
      \texttt{\textasciitilde/.isls/metrics.jsonl} has an entry 
      with \texttt{source: "cli"}.
\item Run a Cockpit forge --- verify entry with 
      \texttt{source: "cockpit"}.
\item \texttt{isls metrics} --- shows summary.
\item \texttt{isls metrics --compare} --- shows comparison table 
      (even with 1 entry per source).
\end{enumerate}


% ──────────────────────────────────────────────────────────────────────────────
\subsection{W5: Studio Frontend}
\label{sec:w5}
% ──────────────────────────────────────────────────────────────────────────────

\subsubsection{Goal}

Upgrade the embedded Studio HTML with a chat interface, readiness 
indicator, forge button, progress view, and result browser.

\subsubsection{Approach}

The Studio is a single embedded HTML file 
(\texttt{static/studio.html}). It uses vanilla JS with no build 
step. The existing Studio has 9 views. D7 adds/modifies 3:

\begin{enumerate}[nosep]
\item \textbf{Architect View (new):} Chat interface on the left. 
      Readiness panel on the right. Entity list below readiness. 
      Forge button at bottom (disabled until ready, green when 
      ready).
\item \textbf{Forge Progress View (new):} Progress bar showing 
      nodes processed. Log stream (one line per node). Compile 
      status indicator. Pulsing animation during LLM calls.
\item \textbf{Result View (modify existing Project view):} File 
      tree of generated project. Click to view file content. 
      Stats summary (files, LOC, tokens, duration).
\end{enumerate}

\subsubsection{Chat UI}

\begin{itemize}[nosep]
\item Message input at bottom (textarea + send button).
\item Messages displayed as chat bubbles (user = right/blue, 
      assistant = left/gray).
\item Entity cards appear inline when the assistant reports 
      new/modified entities.
\item Typing indicator during LLM call.
\item Send button disabled while waiting for response.
\end{itemize}

\subsubsection{Readiness Panel}

\begin{itemize}[nosep]
\item Circular progress indicator (0--100\%).
\item Checklist of criteria (green check / red X / gray pending).
\item Suggestions shown as yellow hints.
\item ``FORGE'' button at bottom: gray when not ready, 
      pulsing green when ready.
\end{itemize}

\subsubsection{Progress View}

\begin{itemize}[nosep]
\item Triggered automatically when forge starts.
\item Progress bar: \texttt{21/50 nodes [=========>          ] 42\%}
\item Each node appears as a log line with icon:
      \begin{itemize}[nosep]
      \item Wrench icon for structural nodes.
      \item Brain icon for LLM nodes.
      \item Checkmark for completed.
      \end{itemize}
\item Compile status: spinning $\to$ green check / red X.
\item Auto-transitions to Result View on completion.
\end{itemize}

\subsubsection{W5 Verification}

\begin{enumerate}[nosep]
\item Open \texttt{http://localhost:8420/studio} in browser.
\item Navigate to Architect view.
\item Type a description, see entity cards appear.
\item See readiness indicator update.
\item Click Forge when ready.
\item See progress view with node-by-node updates.
\item See result view with file browser.
\end{enumerate}


% ══════════════════════════════════════════════════════════════════════════════
\section{File Inventory}
% ══════════════════════════════════════════════════════════════════════════════

\subsection{New Files}

\begin{longtable}{llp{7cm}}
\toprule
\textbf{Crate} & \textbf{File} & \textbf{Purpose} \\
\midrule
isls-gateway & \texttt{src/session.rs}
  & Session state, management, readiness check \\
isls-gateway & \texttt{src/architect.rs}
  & LLM conversation protocol, AppSpec merging \\
isls-forge-llm & \texttt{src/metrics.rs}
  & Generation metrics logging \& comparison \\
isls-cli & \texttt{src/cmd\_metrics.rs}
  & \texttt{isls metrics} CLI subcommand \\
\bottomrule
\end{longtable}

\subsection{Modified Files}

\begin{longtable}{llp{7cm}}
\toprule
\textbf{Crate} & \textbf{File} & \textbf{Change} \\
\midrule
isls-gateway & \texttt{src/lib.rs}
  & Session routes, session store in AppState \\
isls-gateway & \texttt{src/static/studio.html}
  & Architect view, progress view, result view \\
isls-forge-llm & \texttt{src/staged\_closure.rs}
  & Progress callback hook \\
isls-forge-llm & \texttt{src/lib.rs}
  & \texttt{pub mod metrics;} \\
isls-cli & \texttt{src/main.rs}
  & Metrics command, metrics logging in forge paths \\
\bottomrule
\end{longtable}


% ══════════════════════════════════════════════════════════════════════════════
\section{The 12 Rules --- D7 Compliance}
% ══════════════════════════════════════════════════════════════════════════════

\begin{longtable}{clp{9cm}}
\toprule
\textbf{\#} & \textbf{Status} & \textbf{D7 Impact} \\
\midrule
1  & \color{isls-green}\checkmark & LLM surface unchanged.
     The architect LLM call is for conversation, not code gen. \\
2  & \color{isls-green}\checkmark & No module changes in 
     forge-llm (except metrics + callback). \\
3  & \color{isls-green}\checkmark & Entity naming convention 
     enforced by architect protocol (PascalCase). \\
4--9  & \color{isls-green}\checkmark & No changes to generation 
     pipeline. Rules apply to generated output, which is unchanged. \\
10 & \color{isls-green}\checkmark & LLM failure in architect 
     mode $\to$ manual fallback. Forge failure $\to$ error event 
     via WebSocket. \\
11 & \color{isls-green}\checkmark & ProvidedSymbol unchanged. \\
12 & \color{isls-green}\checkmark & Forge prompts unchanged. \\
\bottomrule
\end{longtable}


% ══════════════════════════════════════════════════════════════════════════════
\section{Constraints for Claude Code}
% ══════════════════════════════════════════════════════════════════════════════

\begin{enumerate}
\item \textbf{Do not rewrite} \texttt{isls-gateway/src/lib.rs} or 
      \texttt{studio.html} from scratch. Add to them incrementally.
\item \textbf{Do not modify} the forge pipeline (hdag.rs, 
      structural.rs, static\_files.rs, blueprint.rs) except for 
      adding the progress callback to staged\_closure.rs.
\item \textbf{Do not modify} isls-norms, isls-reader, isls-code-topo, 
      isls-hypercube, isls-types.
\item \textbf{Session state is in-memory} (HashMap in AppState). 
      No database, no file persistence for sessions.
\item \textbf{Metrics persistence} uses 
      \texttt{\textasciitilde/.isls/metrics.jsonl} (append-only).
      Use the same \texttt{dirs\_path()} with USERPROFILE fallback 
      from D4.
\item The Oracle calls in architect mode use the existing 
      \texttt{OpenAiOracle} from isls-forge-llm. No new HTTP 
      client.
\item The Studio HTML remains a single embedded file. No external 
      JS frameworks, no build step.
\item All new code MUST compile with \texttt{cargo build --workspace}.
\item \texttt{cargo test --workspace} --- all existing tests pass 
      (192+), plus new tests.
\item Branch: \texttt{claude/implement-isls-d7-spec}.
\item Commit format: \texttt{D7/W\{n\}: <description>}.
\end{enumerate}


% ══════════════════════════════════════════════════════════════════════════════
\section{Dependency Graph}
% ══════════════════════════════════════════════════════════════════════════════

\begin{verbatim}
W1 (Conversation Sessions)
  |
  v
W2 (Readiness Check)
  |
  v
W3 (Live Forge Integration)
  |
  v
W4 (Generation Metrics)   [parallel-safe with W3]
  |
  v
W5 (Studio Frontend)
\end{verbatim}

W1 $\to$ W2 $\to$ W3 is the critical path. W4 can run in parallel 
with W3 (it modifies different files). W5 depends on all previous 
WPs because it renders their outputs.

Recommended: W1 $\to$ W2 $\to$ W4 $\to$ W3 $\to$ W5.


% ══════════════════════════════════════════════════════════════════════════════
\section{Acceptance Criteria}
% ══════════════════════════════════════════════════════════════════════════════

\begin{longtable}{clp{9cm}}
\toprule
\textbf{\#} & \textbf{Pass} & \textbf{Criterion} \\
\midrule
1 & & \texttt{cargo build --workspace} compiles. \\
2 & & \texttt{cargo test --workspace} passes (192+). \\
3 & & Create session + 3 messages + get readiness via REST. \\
4 & & Forge from session produces compiled application. \\
5 & & WebSocket receives forge progress events. \\
6 & & \texttt{isls metrics} shows generation entries. \\
7 & & Studio Architect view renders in browser. \\
8 & & Studio Forge progress view shows node-by-node updates. \\
9 & & All 6 D4 pipeline apps still compile (backward compat). \\
\bottomrule
\end{longtable}

D7 is passed when criteria 1--4 and 6 are fulfilled.
Criteria 5, 7, 8 depend on browser/WebSocket testing.
Criterion 9 is tested by the pipeline bat.


% ══════════════════════════════════════════════════════════════════════════════
\section{Verification Sequence}
% ══════════════════════════════════════════════════════════════════════════════

\begin{enumerate}
\item \texttt{cargo build --workspace}
\item \texttt{cargo test --workspace}
\item \texttt{isls serve --port 8420 --api-key \$KEY}
\item \texttt{curl -X POST http://localhost:8420/api/session 
      -d '\{"api\_key":"..."\}'} $\to$ session ID
\item \texttt{curl -X POST http://localhost:8420/api/session/\{id\}/message 
      -d '\{"message":"Pet shop app"\}'} $\to$ entities + readiness
\item Repeat with refinement messages.
\item \texttt{curl http://localhost:8420/api/session/\{id\}/readiness}
      $\to$ score + criteria
\item \texttt{curl -X POST http://localhost:8420/api/session/\{id\}/forge}
      $\to$ starts forge
\item Verify output directory contains compiled app.
\item \texttt{isls metrics} $\to$ shows entry with source ``cockpit''.
\item Open \texttt{http://localhost:8420/studio} $\to$ Architect 
      view works.
\item Run pipeline bat $\to$ all 6 D4 apps still compile.
\end{enumerate}

\vfill
\begin{center}
\color{isls-gray}\rule{0.5\textwidth}{0.4pt}\\[0.5em]
{\small End of D7 Specification.\\[0.3em]
Der Architekt entwirft. Der Drucker materialisiert.\\
Die Spirale misst ob Gen~2 besser ist als Gen~1.}
\end{center}

\end{document}
